\section{Concurrent tasks and conflicts}
\label{sec:conflicts}

Papilot runs up to $c = 2$ tasks at once, each against the manuscript as it was
when the task started. Accepted text from one task can therefore overlap a
later proposal from another, and \cref{eq:transform} then returns $\bot$.
This section estimates how often that happens.

\subsection{Uniformly placed edits}

\begin{assumption}
\label{ass:uniform}
Each task proposes edits of a common length $\ell$ whose start positions are
independent and uniformly distributed over a document of length $n \gg \ell$.
\end{assumption}

\begin{proposition}[Conflict probability]
\label{prop:conflict}
Under \cref{ass:uniform}, let two tasks propose $k_1$ and $k_2$ edits. The
number $X$ of incompatible cross pairs has mean
\begin{equation}
  \E[X] = \frac{(2\ell + 1)\, k_1 k_2}{n},
  \label{eq:conflict-mean}
\end{equation}
and the probability that the tasks conflict at all is
\begin{equation}
  \Prob(X > 0)
  = 1 - \Bigl(1 - \frac{2\ell + 1}{n}\Bigr)^{k_1 k_2}
  \approx 1 - \exp\Bigl(-\frac{(2\ell+1)\, k_1 k_2}{n}\Bigr).
  \label{eq:conflict}
\end{equation}
\end{proposition}

\begin{proof}
Two edits of length $\ell$ starting at $s$ and $s'$ violate
\cref{eq:compatible} exactly when $\abs{s - s'} \le \ell$, an event of
probability $(2\ell+1)/n$ up to boundary effects of order $\ell/n$. Summing over
the $k_1 k_2$ cross pairs gives \cref{eq:conflict-mean}. Treating the pairs as
independent gives the product in \cref{eq:conflict}, and
$\log(1-x) = -x + O(x^2)$ gives the approximation, which is the Poisson
approximation $\Prob(X = 0) \approx e^{-\E[X]}$~\cite{feller1968}.
\end{proof}

For a $40$-page paper with $n \approx 1.2 \times 10^{5}$ characters, two tasks
of $k_1 = k_2 = 10$ edits of $\ell = 200$ characters have $\E[X] \approx 0.33$
and conflict with probability about $1 - e^{-0.33} \approx 0.28$.

\subsection{How much overlaps}

When two edits do conflict, the overlap is usually partial.

\begin{proposition}[Expected overlap]
\label{prop:overlap}
Conditioned on a conflict, the offset $d = s' - s$ is uniform on
$\{-\ell, \dots, \ell\}$ and the overlap $\ell - \abs{d}$ has mean
\begin{equation}
  \E\bigl[\ell - \abs{d} \bigm| \abs{d} \le \ell\bigr]
  = \frac{1}{2\ell + 1} \sum_{d=-\ell}^{\ell} \bigl(\ell - \abs{d}\bigr)
  = \frac{\ell^{2}}{2\ell + 1}
  \approx \frac{\ell}{2} .
  \label{eq:overlap}
\end{equation}
\end{proposition}

\begin{proof}
The sum is $\ell + 2\sum_{j=1}^{\ell} (\ell - j) = \ell + \ell(\ell - 1) = \ell^2$.
\end{proof}

A typical conflict therefore involves about half of an edit, which matches the
experience that conflicting proposals usually touch the same sentence rather
than the same word.

\subsection{Clustered edits}

Real revisions are not uniform: an author working on the introduction asks for
several changes there. We model this by letting both tasks place their edits
uniformly within a common region covering a fraction $\nu \in (0, 1]$ of the
document.

\begin{corollary}[Clustering raises conflicts]
\label{cor:clustered}
If both tasks edit only within a region of length $\nu n$,
\begin{equation}
  \Prob(X > 0) \approx 1 - \exp\Bigl(-\frac{(2\ell+1)\, k_1 k_2}{\nu n}\Bigr),
  \label{eq:clustered}
\end{equation}
so that confining two tasks to a tenth of the paper multiplies the expected
number of conflicts by ten.
\end{corollary}

\begin{proof}
Apply \cref{prop:conflict} with $n$ replaced by $\nu n$.
\end{proof}

\Cref{cor:clustered} argues for fewer, longer tasks over many short ones that
touch the same part of a paper, and for finishing the review of one task
before starting another in the same section. It also explains why Papilot
surfaces conflicts from Git rather than resolving them: a resolution would
have to choose between two texts the author has not yet compared.
